\documentclass[11pt,a4paper]{article}

\usepackage[margin=2.5cm]{geometry}
\usepackage[T1]{fontenc}
\usepackage[utf8]{inputenc}
\usepackage{lmodern}
\usepackage{booktabs}
\usepackage{longtable}
\usepackage{tabularx}
\usepackage{array}
\usepackage{xcolor}
\usepackage{listings}
\usepackage{hyperref}
\usepackage{enumitem}

\hypersetup{
    colorlinks=true,
    linkcolor=black,
    urlcolor=blue,
    citecolor=black,
    pdftitle={Test-Driven Development Documentation},
    pdfauthor={Candidate ID: INSERT-ID}
}

\definecolor{codebackground}{RGB}{245,245,245}
\lstset{
    basicstyle=\ttfamily\small,
    backgroundcolor=\color{codebackground},
    frame=single,
    breaklines=true,
    columns=fullflexible
}

\newcolumntype{Y}{>{\raggedright\arraybackslash}X}

\title{Test-Driven Development Documentation\\
\large Secure Software Development Authentication Controls}
\author{Candidate ID: \texttt{INSERT-CANDIDATE-ID}}
\date{Spring 2026}

\begin{document}

\maketitle

\begin{abstract}
This document records the test-driven development (TDD) approach used to verify
and strengthen registration, login, session protection and basic login-attempt
protection in the Secure Software Development Flask application. The document
maps the tests to the authentication and TDD expectations in the MC100 exam,
describes the test environment, and explains the relationship between each
security requirement, test and implementation decision. The focus is deliberately
limited to the local academic application and does not claim production readiness.
\end{abstract}

\tableofcontents
\newpage

\section{Purpose and Exam Alignment}

The MC100 exam requires secure account creation and authentication with modern
password hashing, secure credential verification, and TDD for selected
security-related functionality \cite{mc100exam,mc100guidelines}. In particular,
the work documented here supports:

\begin{itemize}[leftmargin=*]
    \item \textbf{Task A.3.1}: user registration and login with secure password
    hashing.
    \item \textbf{Task A.4.3}: TDD applied to authentication and an additional
    relevant application security scenario.
    \item \textbf{Task A.4.1}: parameterised database access is retained by the
    application models and is covered separately in the penetration-testing work.
\end{itemize}

The tests focus on observable security behaviour rather than implementation
details alone. This makes the evidence useful both for development and for the
exam report.

\section{Scope}

The documented scope is restricted to the following local Flask application
controls:

\begin{itemize}[leftmargin=*]
    \item server-side validation of usernames, email addresses and passwords;
    \item bcrypt password hashing and password verification;
    \item duplicate username and email prevention;
    \item successful and unsuccessful login behaviour;
    \item session-based access control for protected routes;
    \item basic protection against repeated failed login attempts; and
    \item session-cookie settings appropriate for local development.
\end{itemize}

The project remains a local academic system. It does not add JWT, OAuth, MFA,
external identity providers or production deployment infrastructure, because
those are outside the required scope at this stage.

\section{TDD Method and Traceability}

TDD follows the Red-Green-Refactor cycle \cite{beck}:

\begin{enumerate}[leftmargin=*]
    \item \textbf{Red}: write a test that expresses one security requirement and
    observe that the requirement is not yet met.
    \item \textbf{Green}: add the smallest implementation change needed for the
    test to pass.
    \item \textbf{Refactor}: improve clarity while keeping the full test suite
    passing.
\end{enumerate}

The original project already contained baseline registration and login code.
Therefore, this documentation does not claim that all earlier code was
historically written with TDD. Existing behaviour was first protected by
characterisation tests. The new validation and login-attempt controls were then
defined by tests before their implementation was completed. This distinction
keeps the report accurate while demonstrating a real test-first security
workflow for the current iteration.

\section{Test Environment}

The test suite is stored in \texttt{tests/test\_auth.py} and uses \texttt{pytest},
Flask's built-in test client, bcrypt and a temporary SQLite database.

\begin{table}[h]
\centering
\begin{tabularx}{\textwidth}{@{}lY@{}}
\toprule
\textbf{Component} & \textbf{Purpose} \\
\midrule
pytest & Collects and executes the authentication tests. \\
Flask test client & Sends requests without starting a public web server. \\
Temporary SQLite database & Isolates each test from \texttt{users.db} and prevents
test data from affecting the local application database. \\
\texttt{monkeypatch} fixture & Replaces the configured database path with the
temporary database path for each test. \\
\texttt{failed\_login\_attempts.clear()} & Resets the in-memory lockout state so
tests do not influence each other. \\
\bottomrule
\end{tabularx}
\caption{Authentication test environment.}
\end{table}

The isolated database is an important part of repeatable TDD: a failed or
successful test must not depend on a manually created account or on the order in
which a developer previously used the application.

\section{Security Policy Under Test}

\begin{table}[h]
\centering
\begin{tabularx}{\textwidth}{@{}lY@{}}
\toprule
\textbf{Area} & \textbf{Policy} \\
\midrule
Username & 3--30 characters; letters, digits and underscore only. \\
Email & Checked server-side before database access. \\
Password length & Minimum 8 and maximum 16 characters. \\
Password composition & At least one lowercase letter, uppercase letter and digit. \\
Password storage & bcrypt hash using cost factor 12; no plaintext password is stored. \\
Duplicate accounts & Username and email are checked before insertion; the database
also enforces unique email values. \\
Login lockout & Five failed attempts lock the submitted username for 300 seconds. \\
Session cookie & \texttt{HttpOnly=True} and \texttt{SameSite=Lax}. The \texttt{Secure}
flag is false because the assessed system is local and uses HTTP. \\
\bottomrule
\end{tabularx}
\caption{Security policy applied to the local academic application.}
\end{table}

\section{Authentication Test Matrix}

Table \ref{tab:tests} maps each test to its security purpose. The comments inside
the test file state the security requirement, the action under test and the
expected result.

\small
\begin{longtable}{@{}p{0.38\textwidth}p{0.27\textwidth}p{0.27\textwidth}@{}}
\caption{TDD authentication test matrix.}\label{tab:tests}\\
\toprule
\textbf{Test} & \textbf{Security requirement} & \textbf{Expected result} \\
\midrule
\endfirsthead
\toprule
\textbf{Test} & \textbf{Security requirement} & \textbf{Expected result} \\
\midrule
\endhead
\texttt{test\_register\_stores\_bcrypt\_hash\_as\_text} & Passwords are never
stored in plaintext. & A bcrypt text hash is stored and verifies the submitted
password. \\
\texttt{test\_register\_rejects\_invalid\_email} & Email validation is performed
on the server. & No account is created for malformed email input. \\
\texttt{test\_register\_rejects\_invalid\_username} & Usernames have a predictable
safe format. & Whitespace and unsupported characters are rejected. \\
\texttt{test\_register\_rejects\_short\_password} & Passwords meet the minimum
length. & Passwords shorter than 8 characters are rejected. \\
\texttt{test\_register\_rejects\_password\_without\_uppercase\_letter} & Password
composition includes an uppercase character. & Account creation is refused. \\
\texttt{test\_register\_rejects\_password\_without\_lowercase\_letter} & Password
composition includes a lowercase character. & Account creation is refused. \\
\texttt{test\_register\_rejects\_password\_without\_number} & Password composition
includes a digit. & Account creation is refused. \\
\texttt{test\_register\_rejects\_password\_longer\_than\_16\_characters} & Password
length has an upper bound. & Passwords longer than 16 characters are rejected. \\
\texttt{test\_register\_rejects\_duplicate\_username} & Each username identifies
one account. & The duplicate account is not created. \\
\texttt{test\_register\_rejects\_duplicate\_email} & One email cannot be linked
to multiple accounts. & The duplicate account is not created. \\
\texttt{test\_login\_accepts\_valid\_credentials} & Valid credentials authenticate
the user. & The authenticated dashboard is returned. \\
\texttt{test\_login\_rejects\_wrong\_password\_without\_session} & Invalid
passwords do not create sessions. & A generic error is shown and no login session
exists. \\
\texttt{test\_login\_rejects\_unknown\_user\_without\_session} & Unknown users
follow the safe failure path. & A generic error is shown and no login session
exists. \\
\texttt{test\_login\_blocks\_repeated\_failed\_attempts} & Repeated failures are
temporarily rate-limited. & The sixth request is blocked, even with valid
credentials. \\
\texttt{test\_protected\_route\_rejects\_anonymous\_user} & Protected resources
require authentication. & The anonymous client is redirected to login. \\
\texttt{test\_session\_cookie\_configuration} & Session cookies reduce
client-side exposure. & \texttt{HttpOnly} and \texttt{SameSite=Lax} are enabled. \\
\bottomrule
\end{longtable}
\normalsize

\section{Implementation Linked to the Tests}

The following implementation changes were made after the relevant test cases had
been defined:

\begin{itemize}[leftmargin=*]
    \item \texttt{valid\_username()} uses a regular expression to accept only the
    selected 3--30 character username format.
    \item \texttt{valid\_email()} checks that email input has a basic address
    structure before it reaches the database.
    \item \texttt{valid\_password()} enforces the selected length and composition
    policy.
    \item \texttt{hash\_password()} uses \texttt{bcrypt.gensalt(rounds=12)} and
    stores the resulting hash as UTF-8 text.
    \item \texttt{verify\_password()} verifies the submitted password against the
    stored bcrypt hash.
    \item \texttt{login\_is\_locked()}, \texttt{record\_failed\_login()} and
    \texttt{clear\_failed\_login\_attempts()} implement the basic lockout rule.
    \item \texttt{login\_required} continues to protect routes through a Flask
    decorator and redirects anonymous users to the login page.
\end{itemize}

The following excerpt illustrates the password validation rule:

\begin{lstlisting}[language=Python, caption={Password validation rule.}]
def valid_password(password):
    return (
        8 <= len(password) <= 16
        and any(character.islower() for character in password)
        and any(character.isupper() for character in password)
        and any(character.isdigit() for character in password)
    )
\end{lstlisting}

\section{Execution Procedure and Evidence}

From the project root, run the following commands in PowerShell:

\begin{lstlisting}[caption={Test execution commands.}]
python -m venv .venv
.venv\Scripts\Activate.ps1
python -m pip install -r requirements.txt
python -m pytest -q tests/test_auth.py
\end{lstlisting}

The latest observed test execution produced pass markers for all 16 tests. Before
submission, execute the final command locally again and retain a small screenshot
showing the command and final pytest summary. The exam guidance asks for concise
evidence rather than complete screenshots, so one cropped terminal image is
sufficient.

Recommended evidence for the report:

\begin{enumerate}[leftmargin=*]
    \item A small screenshot of the pytest result.
    \item A short code excerpt of one representative test, such as the bcrypt
    storage test or lockout test.
    \item A short code excerpt of \texttt{valid\_password()}.
    \item One or two sentences explaining the Red-Green-Refactor cycle and why
    the selected login lockout scenario matters.
\end{enumerate}

\section{Limitations and Future Work}

The login-attempt protection is intentionally basic. It uses an in-memory Python
dictionary, which is suitable for the local single-process coursework
application. It resets when the application restarts and is not shared across
multiple server instances. A production system would use a persistent or shared
rate-limiting store and would be deployed over HTTPS with a non-hardcoded secret
key and a secure cookie flag.

These production controls are not claimed as part of the current local exam
implementation. Future roadmap stages will add password reset functionality,
secure file handling, further penetration testing and Cerebras API protection.

\section{Conclusion}

The authentication component is covered by a focused suite of 16 repeatable
tests. The suite verifies the most relevant local security properties: validation,
bcrypt hashing, safe login failure behaviour, session protection and a basic
defence against repeated login attempts. The TDD workflow provides traceable
evidence that the current authentication security requirements were specified as
tests and verified after implementation.

\begin{thebibliography}{9}

\bibitem{mc100exam}
Kristiania University College.
\textit{MC100 Secure Software Development Exam Task 2026}.
Course examination paper, Spring 2026.

\bibitem{mc100guidelines}
Kristiania University College.
\textit{MC100 Secure Software Development Grading Guidelines (A--F Scale)}.
Spring 2026.

\bibitem{beck}
Kent Beck.
\textit{Test-Driven Development: By Example}.
Addison-Wesley, 2002.

\bibitem{owasp-authentication}
OWASP Foundation.
\textit{Authentication Cheat Sheet}.
\url{https://cheatsheetseries.owasp.org/cheatsheets/Authentication_Cheat_Sheet.html}

\bibitem{owasp-password-storage}
OWASP Foundation.
\textit{Password Storage Cheat Sheet}.
\url{https://cheatsheetseries.owasp.org/cheatsheets/Password_Storage_Cheat_Sheet.html}

\end{thebibliography}

\end{document}
